% ADAPTIVE-FBDD1
\section*{ADAPTIVE-FBDD1: adaptive read-once attack on the residual-state barrier}

\textbf{STATUS: COMPUTATIONAL finite construction and verification. P versus NP remains OPEN.}

\paragraph{Model.}
An FBDD is a deterministic read-once branching program: the next queried variable may depend on the computation path, but no root-to-sink path queries an input variable more than once. This strictly strengthens the fixed-order OBDD-style model used in the first residual-state experiment.

The target is the exact-weight-eight promise classifier \(D(n,8)\), with output one for UNSAT. Under DN8-FRONTIER1 the positive inputs are exact partitions of the assignment cube by eight falsifying subcubes; for \(n=4\) these are exactly the 272 perfect matchings of \(Q_4\).

\paragraph{Construction.}
For a partial promise state, store the set of remaining clause bits, the number of selected clauses still required, and the subset of positive exact covers still consistent with the history. A state can terminate SAT when no positive completion survives, and can terminate UNSAT when the number of surviving positive completions equals the number of all legal weight-respecting completions. Otherwise a heuristic selects the next clause bit using child-domain and minority-label statistics. Structurally identical decision nodes are hash-consed.

The reproducible constructor is
\texttt{experiments/universal-dag/adaptive\_fbdd\_dnm8.py}
(commit \texttt{107b03c2b6ddf9bba709c053d004c11ad9307c91}).

\paragraph{D(4,8).}
The resulting FBDD has
\[
\boxed{301}
\]
internal nodes and worst-case path depth
\[
\boxed{17}.
\]
Every one of the 32 clause inputs occurs somewhere in the DAG, but no path queries more than 17 variables.

An independent C++ verifier,
\texttt{experiments/universal-dag/verify\_d4\_fbdd.cpp}
(commit \texttt{e0626b0951533df95ef57e34727834737a13f666}),
checks the syntactic read-once condition and exhaustively evaluates all
\[
\binom{32}{8}=10{,}518{,}300
\]
promised inputs. It accepts exactly the 272 perfect matchings and rejects the remaining 10,518,028 inputs. The generated node table has SHA-256
\[
\texttt{434abcf7d1bf7bc5950abeb9672e09e100bd827aa9886d84e826a315b293b8be}.
\]

The synchronous live-node profile by query count is
\[
1,2,4,8,16,29,46,72,89,94,89,63,42,27,17,8,1,
\]
so this particular adaptive machine has peak simultaneous control width 94, or
\[
\log_2 94\approx6.555
\]
bits to identify its live control node if static program storage is treated separately.

The 301-node construction improves the currently recorded 720-node fixed-order upper bound, but it is not claimed minimal.

\paragraph{D(5,8).}
Using the same architecture with a separately fixed heuristic weight vector gives a
\[
\boxed{2248}
\]
node FBDD of worst-case depth
\[
\boxed{28}.
\]
All 4,380 positive exact covers are accepted, and 100,000 deterministic seeded promised inputs were cross-checked against exact-cover membership. The synchronous live-node peak is 517.

\paragraph{Literature audit.}
G\"unther and Drechsler give exact and heuristic FBDD minimization methods and emphasize the huge optimization search space; FBDD minimization is NP-hard and exact algorithms are limited to small functions. Thus the present values are upper bounds, not optima.

Borodin--Razborov--Smolensky give explicit separations showing that read-twice branching programs can be polynomially smaller than read-once programs. Therefore a read-once lower bound, even if strong, cannot be promoted through RECOMP without an additional theorem.

Dinesh--Koroth--Sarma prove that bitwise decomposable projective dimension characterizes deterministic branching-program size within polynomial factors, while ordinary projective dimension can be exponentially smaller. This remains the most relevant literature target for upgrading multisign residual incidence from FBDD cuts to general branching programs.

Beame--Li--Roy--Suciu show that decision-DNNFs can be converted to FBDDs with quasipolynomial blowup in general, illustrating one restricted route by which FBDD lower bounds transfer to richer knowledge-compilation machinery.

\paragraph{Falsification result.}
The fixed-order residual-state measure is not representation-invariant. Adaptive variable choice cuts the tested \(D(4,8)\) representation from 720 recorded fixed-order nodes to 301 FBDD nodes and allows every computation path to terminate within 17 queries. Therefore any load-bearing state invariant must survive adaptive query geometry.

\paragraph{Remaining hostile controls.}
Read-twice/general branching programs remain untested for this classifier. Fixed \(m=8\) also has the independent polynomial pair-conflict circuit and can be streamed in polynomial time while storing only the constant number of selected clause identities. These are mandatory controls against any unrestricted-hardness interpretation.

\paragraph{Next target.}
Attempt a read-twice/RECOMP compression of the 301-node machine. In parallel, formulate a promise-compatible analogue of bitwise decomposable projective dimension or another general-BP nonmergeability certificate. A useful success would be a lower bound that survives repeated queries; a large collapse under read-twice is equally useful because it identifies the exact resource that a pebbling/time-space layer would need to charge.

\begin{thebibliography}{9}
\bibitem{GD02}
W.~G\"unther and R.~Drechsler,
\emph{Minimization of Free BDDs},
Integration 32(1--2):41--59, 2002.
doi:10.1016/S0167-9260(02)00041-X.

\bibitem{BRS93}
A.~Borodin, A.~Razborov, and R.~Smolensky,
\emph{On Lower Bounds for Read-k-Times Branching Programs},
Computational Complexity 3:1--18, 1993.
doi:10.1007/BF01200404.

\bibitem{DKS16}
K.~Dinesh, S.~Koroth, and J.~Sarma,
\emph{Characterization and Lower Bounds for Branching Program Size Using Projective Dimension},
FSTTCS 2016, LIPIcs 65:37:1--37:14.
doi:10.4230/LIPIcs.FSTTCS.2016.37.

\bibitem{BLRS13}
P.~Beame, J.~Li, S.~Roy, and D.~Suciu,
\emph{Lower Bounds for Exact Model Counting and Applications in Probabilistic Databases},
UAI 2013.
\end{thebibliography}

% END_OF_NOTE::ADAPTIVE_FBDD1
